\documentclass[11pt]{article}
\usepackage{amsmath,amssymb}
\usepackage{geometry}
\usepackage{tikz}
\usetikzlibrary{arrows.meta,calc,decorations.markings}
\usepackage{hyperref}
\geometry{margin=1in}

\title{Harmonic Potential Fields for Multi-Robot Exploration:\\
Theory, Covariant Derivatives, and Natural Dispersion}
\author{Working Notes --- W.\ R.\ Kollmyer}
\date{\today}

\begin{document}
\maketitle

\section{The Problem with Classical Potential Fields}

In a standard artificial potential field, each cell $x$ in the grid
carries a scalar potential $\phi(x)$ composed of attractive and
repulsive terms:
\begin{equation}
    \phi(x) = \phi_{\text{attract}}(x) + \phi_{\text{repel}}(x)
\end{equation}
The robot follows the negative gradient $-\nabla\phi$ toward low
potential (unexplored territory) and away from high potential
(obstacles).

This is elegant but has a fatal flaw: \textbf{local minima}.  Where
the attractive gradient from a frontier balances the repulsive
gradient from an obstacle, $\nabla\phi = 0$ and the robot stops
dead.  It hasn't reached its goal; it's trapped in a force
equilibrium.

\section{Harmonic Functions: No Local Minima by Construction}

A function $\phi : \Omega \to \mathbb{R}$ is \textbf{harmonic} on a
domain $\Omega$ if it satisfies the Laplace equation:
\begin{equation}
    \nabla^2 \phi = \frac{\partial^2 \phi}{\partial x^2} +
    \frac{\partial^2 \phi}{\partial y^2} = 0
    \quad \text{in } \Omega
    \label{eq:laplace}
\end{equation}

Harmonic functions have a remarkable property --- the
\textbf{maximum principle}: a harmonic function on a connected domain
attains its extrema \emph{only on the boundary}.  In other words:

\begin{quote}
\emph{There are no local minima or maxima in the interior.}
\end{quote}

Every interior critical point is a saddle, not a trap.  A robot
following $-\nabla\phi$ is mathematically guaranteed to reach a
boundary sink without getting stuck.

\subsection{Why This Works for Exploration}

Set the boundary conditions:
\begin{itemize}
    \item \textbf{Unexplored frontier cells} $\to$ $\phi = 0$
    (attractive sinks)
    \item \textbf{Obstacle / wall cells} $\to$
    $\frac{\partial\phi}{\partial n} = 0$ (Neumann: zero flux, the
    robot slides along walls rather than being attracted or repelled)
    \item \textbf{Explored cells} $\to$ solve Laplace equation
    (\ref{eq:laplace}) to fill in the field
\end{itemize}

The resulting field has smooth gradient flow from every explored cell
toward the nearest frontier, with no local traps, guaranteed.

\subsection{Discrete Laplace Equation on a Grid}

On a grid with cell size $h$, the Laplace equation becomes a
simple averaging relation.  For interior cell $(i,j)$:
\begin{equation}
    \phi_{i,j} = \frac{1}{4}\bigl(\phi_{i+1,j} + \phi_{i-1,j}
    + \phi_{i,j+1} + \phi_{i,j-1}\bigr)
    \label{eq:discrete_laplace}
\end{equation}

Every interior cell's potential is the average of its four neighbors.
This is solved by \textbf{iterative relaxation} (Gauss--Seidel or
Jacobi iteration): initialize $\phi$ with boundary values, then
repeatedly apply (\ref{eq:discrete_laplace}) until convergence.
On a $16 \times 12$ grid (192 cells), this converges in
${\sim}50$--$100$ iterations --- comparable to a single BFS.

The gradient at each cell is:
\begin{align}
    \frac{\partial\phi}{\partial x}\bigg|_{i,j} &\approx
    \frac{\phi_{i+1,j} - \phi_{i-1,j}}{2h} \\
    \frac{\partial\phi}{\partial y}\bigg|_{i,j} &\approx
    \frac{\phi_{i,j+1} - \phi_{i,j-1}}{2h}
\end{align}

The robot at cell $(i,j)$ moves in direction
$-\nabla\phi = \bigl(-\partial_x\phi,\; -\partial_y\phi\bigr)$.

\section{The Multi-Robot Problem: Everyone Chases the Same Sink}

If five robots all follow $-\nabla\phi$ on the same harmonic field,
they all converge to the same nearest frontier and collide.  The
field has no concept of ``this sink is already claimed.''

We need a mechanism that \textbf{naturally disperses} robots across
different frontiers without explicit assignment.

\section{Covariant Derivatives and Curvature-Based Dispersion}

\subsection{Intuition}

The ordinary gradient $\nabla\phi$ tells each robot \emph{which
direction to go}.  The \textbf{second-order structure} of the
field --- how the gradient \emph{changes} from cell to cell ---
tells us something deeper: how many distinct ``channels'' of
flow exist nearby, how sharply they diverge, and whether a robot
is in a bottleneck or an open plain.

If we can measure this second-order structure, we can use it to
spread robots out: send each robot toward a different
\textbf{flow channel} based on the local curvature of the field.

\subsection{The Hessian: Curvature of the Potential}

The Hessian matrix of $\phi$ at a point captures the local curvature:
\begin{equation}
    H(\phi) = \begin{pmatrix}
    \dfrac{\partial^2\phi}{\partial x^2} &
    \dfrac{\partial^2\phi}{\partial x\,\partial y} \\[8pt]
    \dfrac{\partial^2\phi}{\partial x\,\partial y} &
    \dfrac{\partial^2\phi}{\partial y^2}
    \end{pmatrix}
    \label{eq:hessian}
\end{equation}

Since $\phi$ is harmonic, $\nabla^2\phi = 0$, which means:
\begin{equation}
    \frac{\partial^2\phi}{\partial x^2} =
    -\frac{\partial^2\phi}{\partial y^2}
    \label{eq:trace_zero}
\end{equation}

The Hessian is \textbf{traceless}.  Its eigenvalues are
$\lambda$ and $-\lambda$ for some $\lambda \geq 0$.  The
eigenvectors define two orthogonal directions:
\begin{itemize}
    \item The \textbf{positive eigenvector} $\mathbf{e}_+$: the
    direction along which the gradient is \emph{increasing}
    (flow is converging --- a bottleneck or channel)
    \item The \textbf{negative eigenvector} $\mathbf{e}_-$: the
    direction along which the gradient is \emph{decreasing}
    (flow is diverging --- an opening or fork)
\end{itemize}

The magnitude $|\lambda|$ measures \emph{how strongly} the flow
is converging or diverging.

\subsection{Using Curvature to Disperse Robots}

Here is the key idea.  At each cell, decompose the robot's movement
into two components:
\begin{equation}
    \mathbf{v} = \underbrace{-\nabla\phi}_{\text{toward frontier}}
    \;+\; \underbrace{\beta \cdot
    \mathbf{e}_-}_{\text{dispersion along divergent axis}}
    \label{eq:dispersion}
\end{equation}

where $\beta$ is a dispersion gain modulated by:
\begin{enumerate}
    \item \textbf{Peer proximity}: $\beta$ increases when other
    robots are nearby (more need to spread)
    \item \textbf{Curvature magnitude}: $\beta \propto |\lambda|$
    (stronger divergence means more room to spread)
    \item \textbf{Robot identity}: each robot uses a deterministic
    offset based on its ID to pick a slightly different direction
    along $\mathbf{e}_-$, e.g.:
    \begin{equation}
        \beta_k = \beta_0 \cdot |\lambda| \cdot
        \sin\!\bigl(\pi k / N\bigr)
    \end{equation}
    where $k$ is the robot's rank among nearby peers and $N$ is
    the peer count.  Robot 1 drifts left along $\mathbf{e}_-$,
    Robot 2 drifts right, Robot 3 stays centered, etc.
\end{enumerate}

The result: at every fork in the flow field, robots naturally
\textbf{peel off into different channels} proportional to the
curvature, without any explicit assignment or communication
beyond position broadcasts.

\subsection{Connection to Covariant Derivatives}

In differential geometry, the \textbf{covariant derivative}
generalizes the ordinary derivative to curved spaces.  For a
vector field $\mathbf{V}$ (the gradient flow) on a surface
(the potential landscape), the covariant derivative
$\nabla_{\mathbf{u}} \mathbf{V}$ measures how $\mathbf{V}$
changes as you move in direction $\mathbf{u}$:
\begin{equation}
    \nabla_{\mathbf{u}} \mathbf{V} = (\mathbf{u} \cdot \nabla)
    \mathbf{V} + \text{(connection terms from curvature)}
\end{equation}

On a flat grid the connection terms vanish and this reduces to the
directional derivative of the flow field, which is exactly the
Hessian applied to the direction:
\begin{equation}
    \nabla_{\mathbf{u}} (-\nabla\phi) = -H(\phi)\,\mathbf{u}
    \label{eq:covariant}
\end{equation}

Equation~(\ref{eq:covariant}) tells us: \emph{if a robot at
position $x$ moves in direction $\mathbf{u}$, the flow field
rotates by $-H(\phi)\,\mathbf{u}$}.  Large rotation means the
flow is diverging --- there's a fork.  Small rotation means the
flow is uniform --- a corridor.

A robot can compute this quantity locally (the Hessian is just
second-order finite differences on the grid) and use it to decide:
\begin{itemize}
    \item $\|{-H\,\mathbf{u}}\|$ is large $\Rightarrow$ I'm
    at a fork.  Look for a peer-free branch.
    \item $\|{-H\,\mathbf{u}}\|$ is small $\Rightarrow$ I'm
    in a corridor.  Follow the gradient normally.
\end{itemize}

\subsection{Discrete Computation}

On the grid, all quantities are finite differences.  The Hessian
entries at cell $(i,j)$ are:
\begin{align}
    \phi_{xx} &= \frac{\phi_{i+1,j} - 2\phi_{i,j}
    + \phi_{i-1,j}}{h^2} \\
    \phi_{yy} &= \frac{\phi_{i,j+1} - 2\phi_{i,j}
    + \phi_{i,j-1}}{h^2} \\
    \phi_{xy} &= \frac{\phi_{i+1,j+1} - \phi_{i+1,j-1}
    - \phi_{i-1,j+1} + \phi_{i-1,j-1}}{4h^2}
\end{align}

The eigenvalues of the $2 \times 2$ traceless Hessian are:
\begin{equation}
    \lambda = \pm\sqrt{\phi_{xx}^2 + \phi_{xy}^2}
\end{equation}

and the divergent eigenvector (the direction robots should spread
along) is:
\begin{equation}
    \mathbf{e}_- = \frac{1}{\sqrt{\phi_{xx}^2 + \phi_{xy}^2}}
    \begin{pmatrix}
        -\phi_{xy} \\
        \phi_{xx} - \lambda
    \end{pmatrix}
\end{equation}

All of this is $O(1)$ per cell.  No pathfinding, no assignment,
no auction.

\section{Complete Algorithm}

\begin{enumerate}
    \item \textbf{Solve the harmonic field}: set frontier cells to
    $\phi = 0$, walls to Neumann BC, iterate
    (\ref{eq:discrete_laplace}) until convergence.
    \item \textbf{Compute gradient}: $-\nabla\phi$ at each cell
    (central differences).
    \item \textbf{Compute Hessian}: second-order differences,
    eigendecomposition at each cell.
    \item \textbf{Each robot $k$ computes its velocity}:
    \[
        \mathbf{v}_k = -\nabla\phi + \beta_k \cdot |\lambda|
        \cdot \mathbf{e}_- \cdot f(\text{peer proximity})
    \]
    \item \textbf{Move}, mark cells visited, recompute field
    when frontier changes.
\end{enumerate}

The entire pipeline is $O(W \cdot H \cdot I)$ where $W \times H$
is the grid size and $I$ is the relaxation iteration count.  For
$16 \times 12$ with $I = 100$: roughly 200\,000 multiply-adds.
Microseconds on a Pi Zero.

\section{Why This Matters for REIP}

The harmonic field with curvature-based dispersion replaces two
components of the current stack:
\begin{itemize}
    \item \textbf{A* pathfinding} $\to$ gradient following (no
    per-agent search)
    \item \textbf{Leader frontier assignment} $\to$ curvature-based
    dispersion (no explicit allocation)
\end{itemize}

The leader's role reduces to:
\begin{enumerate}
    \item Computing the harmonic field (global map knowledge)
    \item Broadcasting the field (one message, all robots use it)
    \item Trust governance (REIP's actual contribution)
\end{enumerate}

A faulty leader could broadcast a corrupted field, but the
verification is straightforward: each follower can recompute
$\phi$ in its local neighborhood from its own visited-cell set
and check whether the leader's broadcast is consistent.  Field
corruption is detectable the same way bad frontier assignments
are --- by comparing against local knowledge.

\section*{References}

\begin{itemize}
    \item C.~I. Connolly and R.~A. Grupen, ``The applications of
    harmonic functions to robotics,'' \emph{J.\ Robotic Systems},
    vol.~10, no.~7, pp.~931--946, 1993.
    \item J.-O. Kim and P.~K. Khosla, ``Real-time obstacle avoidance
    using harmonic potential functions,'' \emph{IEEE Trans.\ Robot.\
    Autom.}, vol.~8, no.~3, pp.~338--349, 1992.
    \item S.~M. LaValle, \emph{Planning Algorithms}, Cambridge
    University Press, 2006, Ch.~4 (potential functions) and Ch.~8
    (navigation functions).
\end{itemize}

\end{document}
